\section{Related Work}
\textbf{Question answering} 
Question Answering (QA) aims to provide answers to users' questions in natural language~\cite{zhu2021retrieving, pandya2021question}, and most QA systems are composed of information retrieval (IR) and answer extraction (AE)~\cite{mao2021rider, ju2022grape, liu2022heterogeneous}. In IR, the system searches for query-relevant factual passages using heuristic methods (BM25)~\cite{robertson2009probabilistic} or neural-ranking ones (DPR)~\cite{karpukhin2020dense}. In AE, the final answer is extracted usually as a textual span from related passages. Although this framework has been broadly applied in O-QA~\cite{mao2021rider} and D-QA~\cite{xu2020layoutlmv2, mathew2021docvqa}, no previous work focus on MD-QA, which demands alternatively reasoning and retrieving knowledge from multiple documents. To tackle this issue, we construct KGs to encode logical associations among different passages across documents and design an LLM-based graph traversal agent to alternatively generate the reason and visit the most matching passage node.


\noindent\textbf{Pre-train, Prompt, and Predict with LLMs}
With the emergence of LLMs, the paradigm of `pre-train, prompt, predict' has gained magnificent popularity in handling a wide spectrum of tasks~\cite{gururangan2020don, liu2023pre, yu2023large}. This approach begins with pre-training LLMs by pretext tasks to encode world knowledge into tremendous parameters~\cite{wu2023survey} followed by a prompting function to extract pertinent knowledge for downstream tasks~\cite{yang2023harnessing}. Recent advancements explore different prompting strategies to enhance LLMs' reasoning capabilities~\cite{wei2022chain, jin2023large}. In contrast to that, our work offers a novel perspective by transforming the prompt formulation into the KG traversal. 
% To the best of our knowledge, this is the  work exploring the potential of using graph structure in enhancing prompt design for LLMs.



